% Generated draft: requires mathematical and source review.
\documentclass[11pt]{article}
\usepackage[margin=25mm]{geometry}
\usepackage{fontspec}
\setmainfont{Latin Modern Roman}
\usepackage{amsmath,amssymb,mathtools,unicode-math}
\setmathfont{Latin Modern Math}
\usepackage{enumitem,needspace,xurl,hyperref}
\setlength{\parindent}{0pt}
\setlength{\parskip}{5pt}
\setlength{\emergencystretch}{3em}
\newcommand{\N}{\mathbb{N}}
\newcommand{\Z}{\mathbb{Z}}
\newcommand{\Q}{\mathbb{Q}}
\newcommand{\R}{\mathbb{R}}
\newcommand{\C}{\mathbb{C}}
\newcommand{\F}{\mathbb{F}}
\newcommand{\A}{\mathbb{A}}
\newcommand{\PP}{\mathbb{P}}
\newcommand{\E}{\mathbb{E}}
\newcommand{\eps}{\varepsilon}
\newcommand{\dd}{\,\mathrm{d}}
\begin{document}
\section*{UnsolvedMath ID: 30000913}
\textbf{Source problem code:} catalogue entry
\subsection{Definitions and mathematical statement}
Let $X$ be a Davis complex and $G=\operatorname{Aut}(X)$. In the usual setting, a Davis complex is the Coxeter cell complex associated with a Coxeter system, built from its spherical parabolic cosets. The source does not specify that system or additional restrictions on $X$.

To formulate covolumes, equip $G$ with the intended topology of convergence on finite subcomplexes and require the local compactness assumptions needed for Haar measure, as in the locally finite setting. Fix one Haar normalization $\mu$ on $G$. A lattice is a discrete subgroup $\Gamma\leq G$ for which the quotient has finite invariant measure, with covolume written $\mu(\Gamma\backslash G)$. It is uniform if $\Gamma\backslash G$ is compact and nonuniform if the quotient is noncompact. The notation for quotient measure is conventional; the same normalization must be used for all lattices under comparison.

For the given $X$ and $G$, determine whether uniform lattices exist, whether nonuniform lattices exist, and how to construct those that exist. Determine the covolume set
\[
\mathcal V_\mu(G)=
\{\mu(\Gamma\backslash G):\Gamma\leq G\text{ is a lattice}\}
\subset(0,\infty).
\]
In particular, decide whether this set is discrete. In the usual topological meaning this asks whether each $v\in\mathcal V_\mu(G)$ is isolated in that set. Stronger claims of absence of all positive accumulation points or of a positive lower bound should not be silently substituted; the report's intended convention needs confirmation. Changing the fixed Haar normalization rescales the entire set by a common positive factor.

The question is a determination problem as $X$ varies, not a claim that every such $G$ has both kinds of lattices. The class of Davis complexes, the precise automorphism group (for example cellular or type-preserving), its topology, and the hypotheses permitting the quotient invariant measure must be specified for a fully uniform classification.

The supplied literature triage checked 2026-08-21 labels the record partially solved. It attributes constructions of uniform and nonuniform lattices and nondiscreteness phenomena for many Davis complexes to Thomas, while leaving arbitrary cases and full covolume sets unresolved. These are dated attributed progress claims.
\subsection{Short English statement}
For a Davis-complex automorphism group, construct and classify the possible uniform and nonuniform lattices and determine their Haar-normalized covolumes, including discreteness. The source leaves the precise complex and automorphism-group hypotheses unspecified.
\subsection{Sources}
\begin{itemize}
\item Local UnsolvedMath registry, permanent ID 30000913, source code catalogue entry; supplied complete record. Its corrected-verified note reports source checking and normalization without changing the source-level question.
\item Buildings: Interactions with Algebra and Geometry (2008), Oberwolfach Reports, DOI 10.4171/owr/2008/03, pp. 119--172: \url{https://doi.org/10.4171/owr/2008/03}.
\item Anne Thomas: \url{https://arxiv.org/abs/0807.3312}, supplied for the lattice and covolume progress assessment.
\end{itemize}
\end{document}
